\documentclass{css}
%\documentclass[english]{css}

\usepackage[dvips]{graphicx}
\usepackage{latexsym}

\def\|{\verb|}

\newcommand{\cssyear}[0]{2026}
\newcommand{\cssname}[0]{CSS 2026}
\newcommand{\cssversion}[0]{2026/06/01}
\newcommand{\cssemail}[0]{css2026-office@iwsec.org}

\begin{document}

\title{Ghidra High P-codeに基づく関数の意味的等価性判定}
\etitle{Determining Function Semantic Equivalence Using Ghidra High P-Code}

\affiliate{AFFIL}{名古屋工業大学\\
Nagoya Institute of Technology}

\author{斎藤 勇斗}{Hayato Saito}{AFFIL}[h.saito.652@stn.nitech.ac.jp]
\author{掛井 将平}{Shohei Kakei}{AFFIL}
\author{齋藤 彰一}{Shoichi Saito}{AFFIL}

\begin{abstract}
バイナリを対象とするOSS/TPL検出では，対象バイナリの関数と既知ライブラリの関数との類似性から候補を求め，ライブラリ単位に集約する．既存のTPL候補生成手法は大規模な候補探索に適している一方，候補となった2関数が同じ入力に対して同じ出力を返すかは検証しない．そこで本研究は，TPL検出器が出力した関数候補を，集約前に定義した観測範囲で機能的に等価か再検証する．提案する証拠閉包型クロスABI等価性検証は，GhidraのHigh P-codeから出力関係を構成し，異なるABIの物理状態を共通の意味状態へ射影する．経路，関数呼出し，メモリ効果および出力差の局所証拠を同一対象・同一試行へ閉包し，すべての証明義務と試行全体が正常完了した場合だけEQを確定する．NOT\_EQは，指定バージョンの候補から除外する根拠として利用する．現段階ではEqBenchの等価な代表ペアの証拠を追跡し，EQとUNKNOWNを分けた機構から3つの検証仮説を導いた．非等価ペアに対するNOT\_EQ判定能力は今後評価する．
\end{abstract}

\begin{jkeyword}
バイナリ解析，プログラム等価性，High P-code，SMT
\end{jkeyword}

\begin{eabstract}
Binary-level OSS/TPL detectors identify candidate functions by comparing a target binary with known libraries and then aggregate the candidates at the library level. Although existing TPL candidate-generation methods scale to large candidate sets, they do not verify whether two candidate functions produce the same outputs for the same inputs. We therefore propose evidence-closure cross-ABI equivalence checking for verifying TPL function candidates before aggregation. The method builds output relations from Ghidra High P-code and projects ABI-specific physical states to common semantic states. It returns EQ only when local evidence for paths, calls, memory effects, and output equality is closed over the same target and analysis attempt, and the complete analysis run terminates normally. A NOT\_EQ result provides evidence for excluding a candidate for the specified library version. At the current stage, we traced evidence for representative equivalent pairs from EqBench and derived three testable hypotheses from the mechanisms that separated EQ from UNKNOWN. We will evaluate NOT\_EQ decisions on non-equivalent pairs in future work.
\end{eabstract}

\begin{ekeyword}
binary analysis, program equivalence, High P-code, SMT
\end{ekeyword}

\maketitle

\section{はじめに}

ソフトウェアに含まれるOSSおよびThird-Party Library（TPL）をバイナリから特定する技術は，既知脆弱性の影響調査，ライセンス管理およびSoftware Bill of Materials（SBOM）の作成に用いられる．バイナリレベルのTPL検出では，対象バイナリの関数と既知ライブラリの関数との類似性を求め，得られた候補をライブラリやバージョンの単位に集約する．

既存手法は，命令列，制御構造，意味的特徴，関数間関係などを用いて，大量の既知関数から候補を探索する．しかし，類似性が高いことは，候補となった2関数が同じ入力に対して同じ出力を返すことを保証しない．動作の異なる候補が集約に加われば，ライブラリのバージョンを誤って特定する可能性がある．

そこで本研究は，TPL検出器が出力した関数候補を集約前に再検証する．GhidraのHigh P-codeから，既定の観測契約が可視とする戻り値，メモリ，入出力および終了に関係する式と効果を構成し，同じ論理入力に対する差の有無をSMTソルバで調べる．安全な出力差を確認できた候補はNOT\_EQとして指定バージョンの候補から除外する．差が見つからない場合でも，観測対象の経路，関数呼出しまたは副作用を十分に表現できなければEQとせずUNKNOWNとする．

評価では，EqBenchのCコードをARM64とMIPS64，ClangとGCC，O0--O3の16条件でビルドした．等価な4関数の全480ペアについて，関数全体の特徴，最適化，コンパイラおよびISAを分けて集計し，EQとUNKNOWNが分かれる要因を調べた．実High P-codeと未完了証拠を用い，ループ，戻り値格納域および関数呼出しが判定に与える影響を分析する．

\section{背景と関連研究}

\subsection{バイナリレベルのTPL検出}
バイナリレベルのTPL検出は，対象バイナリから文字列，命令，関数の意味的特徴，関数呼出し関係などを抽出し，既知ライブラリとの対応を求める．ModX\cite{modx}は意味的特徴とプログラムのモジュール化を組み合わせて部分的に取り込まれたTPLを検出し，BinaryAI\cite{binaryai}は関数の意味的表現とリンク時の局所性を用いて候補を絞り込む．これらはコンパイラや最適化による表現の違いを越えて候補を探すことを目的としており，候補となる関数対をどれだけ高精度に見つけ，ライブラリ単位へ適切に集約するかに主眼を置く．一方，候補となった関数対が実際に同じ入力に対して同じ出力を返すかどうかは，特徴や構造の類似度からは直接確認できない．本研究は，これらの手法が出力する候補を対象に，実行結果の一致を後段で再検証する．

\subsection{バイナリの意味比較}
BinHunt\cite{binhunt}は基本ブロックをシンボリック実行し，意味的に対応するブロックと制御フローから関数の類似性を評価する手法であり，iBinHunt\cite{ibinhunt}はこれを関数間制御フローの差分探索へ拡張した．これらはバイナリ全体または関数全体を対象に意味的な差分を探索するものであり，既に候補として与えられた関数対を検証する後段処理としては設計されていない．

BinGo\cite{bingo}はクロスアーキテクチャでの類似関数検索を扱い，BLEX\cite{blex}とBinSim\cite{binsim}は実行またはトレースに基づく意味比較を行う．これらは動的解析を前提とするため，実行環境の準備や網羅的な入力生成が必要になる．本研究は静的解析であるHigh P-codeとSMTソルバを用いることで，実行環境なしに入出力の対応関係を検証する．

DDEC\cite{ddec}は，x86の2つのループが同じ動作をするかを検証する手法であり，実行結果から対応関係の候補を求め，SMTで検証する．低水準コード等価性検証\cite{lowlevel-eq}は，低水準コードの意味的一致を扱う．これらはループを含むプログラムの等価性判定という点で本研究と関係するが，対象は個別のループやプログラム断片であり，TPL検出の後段処理として大量の関数対応候補を検証する枠組みではない．

High P-codeの形式意味論は既存研究で検討されている\cite{formal-pcode}．これはHigh P-codeの意味を定義するものであり，関数間の等価性判定手法ではないが，本研究がHigh P-codeを解析対象とする際の意味論的な裏付けとなる．

\subsection{既存研究の課題と本研究の位置づけ}
バイナリレベルのTPL検出研究は，候補生成とライブラリ単位への集約の精度向上に主眼を置いており，候補となった関数対が実行結果の面でも一致するかどうかは十分に検証していない．一方，バイナリの意味比較に関する研究は，関数やループ単位の等価性判定手法を提供するが，TPL検出で得られる大量の候補を対象とした後段検証としての位置づけでは設計されていない．

本研究は，TPL検出等で得られた関数対応候補を入力とし，対象の2バイナリについて対応する入力に対する観測結果を比較して，候補をライブラリ単位へ集約する前にEQまたはNOT\_EQを確認する後段検証器を提案する点で，これらの既存研究と異なる．判定に必要な情報を確認できない場合はUNKNOWNとし，制限時間内に確定できない場合はTIMEOUTとする．

\section{提案手法}

\subsection{入力と比較範囲}
提案手法は，TPL検出器または関数類似度解析器が出力した関数の組を入力とする静的検証器である．実運用で入力するのは，調査対象のバイナリと，バージョンが既知のTPLから作成した参照バイナリであり，調査対象のソースコードは使用しない．命令列の一致ではなく，ABI，Ghidraが復元した型と格納幅，およびHigh P-codeから得られる振る舞いを比較する．

関数$f$について，2つのバイナリで対応づけた共通の入力範囲を$D_f$とする．各バイナリの引数レジスタやスタック領域を同じ引数へ対応づけ，ABIとGhidraが復元した型・格納幅に基づいて，High P-code上の戻り値格納場所から比較する値を取り出す．比較する振る舞い（以下，観測対象）は，戻り値，仮引数から到達可能なメモリ，大域領域，スレッド局所領域（TLS），ヒープ，\texttt{volatile}アクセス，入出力，ユーザ定義P-code操作，正常・異常終了およびトラップ（trap）とする．\texttt{errno}と浮動小数点環境は既定では対象外であり，並行実行は範囲外である．値，対象領域または効果を比較式で表せない場合は「観測対象外」とせずUNKNOWNにする．入力$x\in D_f$に対する観測結果を$\widehat{O}_i(x)$とすると，本稿では
\[
  \forall x\in D_f,\quad \widehat{O}_1(x)=\widehat{O}_2(x)
\]
を確認できた場合に，2つのバイナリの対象関数を，本稿で定めた観測範囲において等価とする．

異なるABI間では，引数の順序，型，符号およびビット幅を照合する．拡張や切詰めを伴う場合は，変換前後の値の対応と，比較対象としない上位ビットが結果に影響しないことを確認する．入力，戻り値または副作用を同じ対象として対応づけられない場合はUNKNOWNとする．

\subsection{High P-codeからの式抽出と比較}
各関数をGhidraで逆コンパイルし，\texttt{HighFunction}，基本ブロック，\texttt{PcodeOpAST}および\texttt{Varnode}から制御フローとデータ依存を取得する．逆コンパイル結果やP-codeの文字列表現は，演算の意味を決める情報として用いない．RETURNから値の定義を逆向きにたどり，整数演算，比較，分岐条件および\texttt{MULTIEQUAL}を式へ変換する．出力対象のメモリ更新は，2つのバイナリで同じメモリ領域として対応づけた対象と，その書込み範囲を特定できる場合に扱う．

外部関数呼出しは，呼出し先，引数，戻り値および出力に影響するメモリ更新を2つのバイナリの対象関数間で対応づけ，その関係を比較式へ組み込む．メモリは，比較対象とするオブジェクトと書込み順序をHigh P-codeから取得して対応づける．ループは，反復全体の入力と出力の関係を式で表せる場合に比較へ組み込む．CALL，メモリ更新，入出力，終了，トラップ，循環する\texttt{MULTIEQUAL}と反復回数，または演算の関係を確定できない場合はUNKNOWNとする．

2つのバイナリの対象関数から得た終了経路の組合せごとに，経路条件が同時に成立するかを調べる．同時に成立しない組合せを除き，成立し得る組合せについて観測出力が異なる条件をSMTソルバへ与える．EQを出す前には，\tabref{tab:eq-conditions}の条件を確認する．これにより，個々の出力差問い合わせがUNSATでも，CFG上の実行可能性を確認できない経路や，比較式に含まれていない副作用が残る場合をUNKNOWNにできる．解析結果と設定も記録し，異なる実行の結果が混在しないようにする．

\begin{table}[t]
\centering
\caption{EQを判定する前に確認する条件}
\ecaption{Conditions Checked before Returning EQ.}
\label{tab:eq-conditions}
\scriptsize
\setlength{\tabcolsep}{3pt}
\begin{tabular}{@{}p{19mm}p{54mm}@{}}
\hline\hline
条件 & 確認内容 \\
\hline
入力対応 & 引数の意味，型，幅および符号を対応づけられる \\
終了・経路 & 対象とする終了経路を解析できる \\
値・副作用 & 比較対象の戻り値，メモリおよび外部効果を表現できる \\
入力範囲の網羅 & 終了経路の条件が入力領域を覆う \\
出力差 & 到達可能な全経路組で不一致条件がUNSATとなる \\
解析結果の対応 & 比較する結果が，記録された同じ関数対，設定および解析実行に対応する \\
\hline
\end{tabular}
\end{table}

\subsection{判定結果とTPL候補への適用}
判定結果は次の4種類である．定義は本節に集約する．
\begin{description}
  \item[EQ] \tabref{tab:eq-conditions}の全条件を確認し，2つのバイナリで対応づけた同じ入力に対する観測出力の不一致条件がすべてUNSATである．
  \item[NOT\_EQ] 到達可能な入力で観測出力が異なり，その差が比較式に含まれていない値，経路または副作用に依存しない．
  \item[UNKNOWN] EQまたはNOT\_EQの確定に必要な入力対応，経路，値，呼出しまたは副作用を表現または確認できない．
  \item[TIMEOUT] ペア全体の制限時間内に，EQ，NOT\_EQまたはUNKNOWNのいずれかを有効な最終判定として確定できない．
\end{description}
個別のSMT問い合わせが制限時間を超えても，後続の解析によって有効な最終判定を得られた場合はTIMEOUTとしない．対象バイナリの読込み，GhidraまたはSMTソルバの起動，解析中の処理，もしくは判定結果の保存・整合性確認に失敗し，EQ，NOT\_EQ，UNKNOWNまたはTIMEOUTのいずれも有効な最終結果として記録できなかった場合は，ERRORとして別に記録する．NOT\_EQは指定バージョンの候補を除外する根拠となるが，同じTPLの別バージョンまで否定しない．UNKNOWNとTIMEOUTは候補を維持し，後続処理へ理由を渡す．

\section{評価と考察}

\subsection{評価課題と実験条件}
本稿では，提案手法単体の関数レベルの判定能力を評価する．実際のTPL検出器が生成した候補集合や，ライブラリ単位の検出精度向上は評価対象としない．次のResearch Question（RQ）に沿って評価する．
\begin{description}
  \item[RQ1] 等価ペアと非等価ペアを，それぞれどの範囲で安全にEQまたはNOT\_EQと判定できるか．
  \item[RQ2] 制御構造，状態変化およびビルド差と判定結果には，どのような関連があるか．
  \item[RQ3] UNKNOWNとTIMEOUTは，どの証拠不足または計算上の限界から生じるか．
  \item[RQ4] 関数レベルの結果は，指定バージョンとの厳密な対応候補を選別する後段検証へどのように利用できるか．
\end{description}

\paragraph{等価ペアの評価}
等価セットでは，同一意味を持つ基準ソースを16条件でビルドし，生成されたバイナリ間のEQ判定能力を評価する．EQBENCH上の基準ソース名は\texttt{new.c}または\texttt{newV.c}である．ビルド条件は，命令セット2種（ARM64，MIPS64），コンパイラ2種（Clang，GCC），最適化4段階（O0--O3）の直積であり，1ソース当たり16バイナリを生成する．1関数について異なる2バイナリの全組合せを比較するため，比較数は$16 \times 15 / 2 = 120$ペアである．現在のEqログは，この120ペアの結果を原関数ごとに記録したものである．

同一ソース由来であることだけを等価性の根拠とはせず，未定義動作，処理系定義動作，整数・ポインタ幅および符号付きオーバーフローへの依存を確認し，比較する入力領域と観測結果の幅を明示する．

\paragraph{非等価ペアの評価}
非等価セットでは，意味を変更する前後のソースを新側・旧側として比較し，NOT\_EQ判定能力を評価する．EQBENCH上では，\texttt{new.c}と\texttt{old.c}，または\texttt{newV.c}と\texttt{oldV.c}の組を用いる．定数・演算子，条件境界，分岐出力，ループ更新，メモリ更新または呼出し効果など，導入された意味差の種類も記録する．Neq実験は現在検証中であり，組合せ方法と最終比較数が確定するまでは，NOT\_EQ判定能力の結果を確定値として記載しない．

% TODO(neq-design): 新側16種と旧側16種の全組合せを比較するか，同一ビルド条件だけを対応させるか，new--oldとnewV--oldVを別集計するか，および1関数当たりの最終比較数を，実験完了後に確定する．

\paragraph{共通の評価単位}
EqとNeqは母数と組合せ規則が異なるため別集計とし，全実験を合算した指標は補助情報に限定する．各評価系では，ペア加重集計とベンチマークケース平均を併記する．ペア加重集計は全比較ペアを数える．ベンチマークケース平均では，Eqは基準ソース中の1関数を1ケース，Neqは意味変更前後の関数対を1ケースとし，各ケースの判定率を計算してから平均する．ビルド差の影響は同一ケース内で比較し，関数または変更ケース自体の難しさとコンパイラ・命令セット差を分離する．ERROR，未完了および成果物欠損は実験健全性として別に集計し，対象から黙って除外しない．

\paragraph{実行条件}
EqとNeqに共通して，解析プログラムの版，実行環境，並列数，ペア当たりの制限時間およびソルバ設定を固定する．

% TODO(experiment): EQBENCHの版と等価性の定義，未定義・処理系依存動作の除外条件，入力領域・出力幅の対応，対象関数数，解析プログラムの版，並列数，制限時間，実行環境および再現手順を，最終成果物の識別情報とともに固定する．

\subsection{判定能力（RQ1）}
4値判定を単一のaccuracyへまとめず，等価ground truthと非等価ground truthに分けて示す．等価セットでは，EQ，NOT\_EQ，UNKNOWN，TIMEOUTの件数と割合を示し，EQ率と判定網羅率を主要指標とする．等価ペアをNOT\_EQとした結果は危険な誤判定として独立に報告する．非等価セットでも4値の件数と割合を別表で示し，NOT\_EQ率と判定網羅率を主要指標とする．非等価ペアをEQとした結果は危険な誤判定である．判定網羅率は，各セットにおける「（EQ件数＋NOT\_EQ件数）／対象ペア数」とする．

本評価では，解析が生成した入力割当てを対象バイナリへ与える独立実行は行わない．したがって，解析内部で差の安全性を確認した結果を「反例確認率」として重ねて数えず，NOT\_EQに対して入力割当てと観測差の根拠を追跡可能な成果物として保存できた割合を，反例成果物生成率として補助的に示す．これはNOT\_EQの正しさを別経路で検証する指標ではなく，判定根拠の監査可能性を表す．

UNKNOWNおよびTIMEOUTは，誤ったEQまたはNOT\_EQと同じ失敗として扱わない．差が存在する非等価ペアをUNKNOWNとした場合は，非等価性を否定したのではなく，安全な反例として確定するための証拠が不足したものとして分類する．EqとNeqを合算した判定網羅率は，ペア数の多い評価系に支配されるため主要な結論には用いない．危険な誤判定数と判定網羅率を並べることで，保守性と適用範囲の両方を評価する．

% TODO(result-eq): 完了したEq成果物から，ペア加重集計とベンチマークケース平均による4値判定，EQ率，判定網羅率および危険なNOT\_EQ件数を作成する．
% TODO(result-neq): Neq実験完了後に，ペア加重集計とベンチマークケース平均による4値判定，NOT\_EQ率，判定網羅率，反例成果物生成率および危険なEQ件数を作成する．

\subsection{判定を分ける特徴と保留原因（RQ2・RQ3）}
制御構造は，直線的な処理，条件分岐，複数の終了経路，ループおよび再帰に分ける．状態変化は，戻り値のみ，ポインタ出力，メモリ更新および外部関数呼出しに分ける．ビルド差は，コンパイラ，最適化，命令セットおよび値の幅に分け，分岐数，論理式の規模および反復関係の有無も記録する．これらと判定結果のクロス集計はEqとNeqを分けて示す．Eq側のUNKNOWNは等価性を証明するための証拠不足，Neq側のUNKNOWNは観測差を安全な反例として確定するための証拠不足を表すためである．EqとNeqの合算値は補助情報に限定する．クロス集計からは「関連が観測された」と述べるにとどめ，原因は代表例の解析と同一ベンチマークケース内の比較によって説明する．

UNKNOWNについては，最終判定がUNKNOWNとなった比較数を「UNKNOWNペア数」，各ペアに記録された証拠不足の延べ数を「原因出現数」，同一ペア内で同じ分類に属する原因の重複を除いたペア数を「根本原因別ペア数」として区別する．原因は第3節の5分類へ対応づける．1ペアが複数分類の原因を持つことを許す複数回答形式であるため，原因ごとの件数をUNKNOWNペア数で割った割合の合計は100\%になるとは限らない．実装内部の詳細な識別名の延べ数は主結果とせず，5分類への対応を再現するための補助情報とする．

終了経路が増える場合は経路条件の網羅，ポインタ出力や呼出しを含む場合は副作用の把握，異なる命令セット間では値幅と呼出し規約の対応，ループ・再帰では反復関係の証明が判定の境界となり得る．必要な振る舞いを論理式へ変換できない場合はUNKNOWN，式を構成できても制限時間内に解けない場合はTIMEOUTである．結果確定後は，これらの一般論を観測結果と代表例へ統合し，重複する説明を削除する．

% TODO(analysis): EqとNeqを分け，各特徴について，ベンチマークケース，ビルド条件，観測結果，不足した証拠または計算量を同じ順序で示す．UNKNOWNはペア数，原因出現数，根本原因別ペア数を区別し，TIMEOUTは設定時間と処理時間の分布を併記する．

\subsection{後段検証への適用可能性と妥当性への脅威（RQ4）}
関数レベル評価で直接示せるのは，Eqセットでどの範囲の等価ペアを安全にEQとできるか，およびNeqセットでどの特徴を持つ非等価ペアを安全にNOT\_EQとできるかである．NOT\_EQは，当該バージョンの関数とは観測結果が異なることを示すため，正確なバージョン同定を必要とするSBOM作成に有用な負の証拠となる．一方，同じTPLの別バージョンに由来する可能性や，対象バイナリにTPL自体が存在する可能性は否定しない．EQも観測モデルの範囲で関数対応を支持する結果であり，単独でライブラリの同一性を証明するものではない．

したがって，本評価は，TPL検出システム全体のprecisionまたはrecallが向上したことを裏付けない．ライブラリ判定の偽陽性低減を定量化するには，実際のTPL検出器が生成した候補集合へ接続し，関数候補の除外前後で，バージョンを含むライブラリ単位のground truthと比較する統合評価が必要である．

妥当性への脅威として，EQBENCHが実際のTPL候補に現れる関数分布を代表するとは限らないこと，同一ソースでも未定義・処理系依存動作により意味が変わり得ること，観測モデル外の副作用があることを挙げる．現在の主要課題は，終了経路の網羅，呼出し先・メモリ効果の要約，一般的なループ・再帰の反復関係，異なる値幅の対応，反例の安全性確認および論理式の計算量である．これらはUNKNOWNまたはTIMEOUTの原因と対応づけ，関数レベル評価の適用限界として示す．

\section{おわりに}

本稿では，TPL検出器が生成した関数候補を集約前に再検証する手法を示した．GhidraのHigh P-codeから既定契約上の戻り値と効果を構成し，SMTソルバで差を調べる．観測対象の経路，CALL，メモリ，入出力，終了またはループを十分に表現できない場合は，差が見つからなくてもEQとせずUNKNOWNとする．

EqBenchの4関数・480ペアでは，257ペアがEQ，216ペアがUNKNOWN，7ペアがTIMEOUTであった．最適化，コンパイラ，ISAを個別に層別したうえで，4関数全体をループ・CALLの有無でも集計した．解析版A内では，\texttt{ltfive/lib}のUNKNOWNに戻り値観測契約・幅の由来不足が，\texttt{ltfive/client}のUNKNOWNにCALL位置対応と経路被覆・反例安全性が多く現れた．解析版Bの\texttt{LoopMult10/foo}では，非循環45ペアがすべてEQで，循環を含む75ペアがすべてUNKNOWNとなり，ループ要約と反例安全性の未完了が対応した．これらはペア内で重複する原因の該当数であり，版AとBをまたぐ効果量の比較ではない．また，7件のTIMEOUTは\texttt{ltfive/client}に集中したが，保存ログから処理段階やCALLとの因果は特定できない．非等価ペアと実際のTPL候補への適用は今後の課題である．


\begin{thebibliography}{99}
\bibitem{modx}
Yang, C., Xu, Z., Chen, H., Liu, Y., Gong, X. and Liu, B.:
ModX: Binary Level Partially Imported Third-Party Library Detection via Program Modularization and Semantic Matching,
Proc. ICSE, pp.1393--1405, DOI: 10.1145/3510003.3510627 (2022).

\bibitem{binaryai}
Jiang, L., An, J., Huang, H., Tang, Q., Nie, S., Wu, S. and Zhang, Y.:
BinaryAI: Binary Software Composition Analysis via Intelligent Binary Source Code Matching,
Proc. ICSE, Article 224, pp.1--13, DOI: 10.1145/3597503.3639100 (2024).

\bibitem{binhunt}
Gao, D., Reiter, M. K. and Song, D.:
BinHunt: Automatically Finding Semantic Differences in Binary Programs,
Proc. ICICS, LNCS 5308, pp.238--255 (2008).

\bibitem{ibinhunt}
Ming, J., Pan, M. and Gao, D.:
iBinHunt: Binary Hunting with Inter-Procedural Control Flow,
Proc. ICISC, LNCS 7839, pp.92--109 (2013).

\bibitem{bingo}
Chandramohan, M., Xue, Y., Xu, Z., Liu, Y., Cho, C. Y. and Tan, H. B. K.:
BinGo: Cross-Architecture Cross-OS Binary Search,
Proc. FSE, pp.678--689 (2016).

\bibitem{blex}
Egele, M., Woo, M., Chapman, P. and Brumley, D.:
Blanket Execution: Dynamic Similarity Testing for Program Binaries and Components,
Proc. USENIX Security, pp.303--317 (2014).

\bibitem{binsim}
Ming, J., Xu, D., Jiang, Y. and Wu, D.:
BinSim: Trace-Based Semantic Binary Diffing via System Call Sliced Segment Equivalence Checking,
Proc. USENIX Security, pp.253--270 (2017).

\bibitem{ddec}
Sharma, R., Schkufza, E., Churchill, B. and Aiken, A.:
Data-Driven Equivalence Checking,
Proc. OOPSLA, pp.391--406 (2013).

\bibitem{formal-pcode}
Naus, N., Verbeek, F., Walker, D. and Ravindran, B.:
A Formal Semantics for P-Code,
Proc. VSTTE, LNCS 13800, pp.111--128 (2023).

\bibitem{proof-cert-smt}
Mebsout, A. and Tinelli, C.:
Proof Certificates for SMT-based Model Checkers for Infinite-State Systems,
Proc. FMCAD, pp.117--124 (2016).

\bibitem{lowlevel-eq}
Putro, P. A.:
Low-Level Code Semantic Equivalence Checker,
Programming and Computer Software, No.6, pp.53--71,
DOI: 10.7868/S3034584725060047 (2025).
\end{thebibliography}

\end{document}
